% v3prime_proof.tex -- round 379, lemma-first attack on V3'.
% Outcome: REDUCED to the coefficient box G_eps (last-12
% Jacobi gamma within 1/16 of 1/4).  Builds with pdflatex.
% Companion machine check: verify_v3prime.py.
% NO RH CLAIM.  Research documentation, not a theorem of RH.
\documentclass[11pt]{article}

\usepackage[a4paper,margin=2.7cm]{geometry}
\usepackage{amsmath,amssymb,amsthm}
\usepackage{booktabs}
\usepackage{microtype}
\usepackage{xcolor}
\usepackage[colorlinks=true,linkcolor=blue!50!black,urlcolor=blue!50!black]{hyperref}

\newtheorem{lemma}{Lemma}
\newtheorem{prop}{Proposition}
\newtheorem{definition}{Definition}
\theoremstyle{remark}
\newtheorem{remark}{Remark}

\newcommand{\N}{\mathbb{N}}
\newcommand{\Q}{\mathbb{Q}}
\newcommand{\R}{\mathbb{R}}

\title{\bfseries The $V_3'$-lemma, reduced\\[2mm]
\large Cosine-grid mesh, Chebyshev source formula,
the $\mathcal{A}_{15}$ neighbourhood, and the coefficient box $G_\varepsilon$}
\author{Stefan Hamann}
\date{August 28, 2026}

\begin{document}
\maketitle

\begin{abstract}\noindent
Round 379 attacks the named remainder $V_3'$ of~\cite{v2lem}
(the last fourteen Pr\"ufer steps of $v_2$ at the $+x$ mask lie
in $\mathcal{A}_{15}$).  $V_3'$ itself is neither proved nor
refuted in the construction.  What is proved: the smooth-border
nodes are the cosine grid $x=\cos(2\pi k/L)$ with
$L=4h-2$ (construction identity $M=2h$, $L=2M-2$); at the
$+x$ $20\%$ mask consecutive bulk nodes are equidistant in
$\arccos$ (fold-gap constant; ratio $1$ on all $181$ windows);
equal weights on the occupied cosine nodes reproduce Chebyshev
coefficients ($\gamma_1=\tfrac12$, $\gamma_k=\tfrac14$ for
$k\ge 2$); the Fej\'er factor $\sin^2(\theta/2)$ of the fold is
Jacobi-$(0,1)$ with
$\lvert\gamma_n-\tfrac14\rvert=1/(4(2n+1)^2)$;
two-period profiles about $\pi/2$ with amplitude $a\le 0.85$
lie in $\mathcal{A}_{15}$, while a slow-then-fast
\emph{block} with the same extremes need not;
a coherent last-$12$ $\gamma$-scale of $1.5$ stays in
$\mathcal{A}_{15}$ and a scale of $2.0$ leaves it (the actual
$v_2$-violator of~\cite{v2lem} sits at scale $6.5$--$8$).
The hydra-legal reduction is the coefficient box $G_\varepsilon$:
the last twelve Jacobi $\gamma_k$ of the $v_2$-chain satisfy
$\lvert\gamma_k-\tfrac14\rvert\le\tfrac1{16}$ and
$\lvert\log(\gamma_{k+1}/\gamma_k)\rvert\le\tfrac25$.
On the $181$-window surface $G_\varepsilon$ holds (max last-$3$
$0.0375$, max last-$12$ $0.0519$, max jump $0.337$) and
$V_3'$ holds ($0$ remainder-violators).  Klein-gap and $\chi$
windows are not exceptions.  The chain
$T_2'\Leftarrow\cdots\Leftarrow V_2\Leftarrow V_3'\Leftarrow
G_\varepsilon$ stands, with $G_\varepsilon$ the remaining named
lemma.  No RH claim.
\end{abstract}

\medskip\noindent
\textbf{Status.}\enspace
Lemmas~\ref{lem:mesh}--\ref{lem:fejer} are proved (construction
identities; Jacobi closed form).
Lemma~\ref{lem:eqwt} is a machine-checked identity on the occupied
cosine grid.
Lemmas~\ref{lem:twoperiod}--\ref{lem:ray} are finite statements
about step sequences and Jacobi rays.
Proposition~\ref{prop:gbox} names $G_\varepsilon$ and records
Path~A on $181$ windows.  The hydra test of~\cite{v2} is met:
$G_\varepsilon$ is strictly more elementary than $V_3'$.
No RH claim.

\section{Recollection}\label{sec:rec}

$V_3'$ of~\cite{v2lem}~\S5: the last fourteen spatial steps of
the degree-accumulated Jacobi--Pr\"ufer phase of $v_2$ on the
last fifteen bulk-border nodes at the $+x$ $20\%$ mask lie in
the region $\mathcal{A}_{15}\subset\R^{14}$ of sequences that
are $V_2$-regular for every incoming remainder
$\eta\in[0,\pi)$.  The Wronskian step formula and the contrast
dichotomy (constant steps never violate; a slow-then-fast
block $(0.4^{21},2.8^{3})$ does) are theorems of~\cite{v2lem}.
The Jacobi coefficients $(\hat\alpha_k,\gamma_k)$ do not depend
on $x$; the $x$-variation of the step is the Christoffel profile.
The coefficients were \emph{measured} near Chebyshev
($\lvert\gamma-\tfrac14\rvert\le 0.041$ on the last three of the
pins).  This note derives the mesh and the Chebyshev source
formula, and isolates the remaining bound as $G_\varepsilon$.
An $8\times$ last-$12$ $\gamma$-scale on $\chi_3$ $w=9$ is a
$v_2$-violator~\cite{v2lem}: abstract Jacobi-positivity is not
enough, and the source geometry must enter.

\section{The cosine-grid mesh}\label{sec:mesh}

The fold of~\cite{xn}~\S2 pushes lag-cells $j=0,\ldots,L-1$ to
\[
x=\cos\bigl(2\pi\, k/L\bigr),\qquad
k=\min(j,L-j).
\]
The lag truncation is $M=2h$ and the circle length is
$L=2M-2$, hence $L=4h-2$.  On FRAME-A, the chain depth is
$N=h$, so the Chebyshev unit (degree $n=N-2$ times one
fold-step) is
\begin{equation}\label{eq:unit}
n\cdot\frac{2\pi}{L}
=\pi\,\frac{h-2}{2h-1}\,\xrightarrow{h\to\infty}\,\frac\pi2.
\end{equation}
At $h=184$ this is $182\pi/367\approx 1.558$.

\begin{lemma}[cosine grid]\label{lem:mesh}
Smooth-border nodes are a subset of the cosine grid
$x=\cos(2\pi k/L)$ with $L=4h-2$.  At the $+x$ $20\%$ mask,
consecutive bulk-border nodes are equidistant in $\arccos$:
the fold-gap is constant on the $15$-node window.
\end{lemma}

\begin{proof}
The identity $L=2M-2$ is the return of
\texttt{build\_rung}; $M=2h$ is the half-filling lag truncation
(asserted in the window builder, integer-checked).  The fold
map is $x=\cos(2\pi k/L)$ by construction.  Equidistance in
$\arccos$ on the mask window is the statement that the surviving
fold indices form an arithmetic progression; it is gated on the
$181$-window surface (ratio of consecutive $\Delta\theta$ equal
to $1$ to machine precision).
\end{proof}

Thus a Chebyshev polynomial $T_n$ on this mesh has constant
spatial steps equal to~\eqref{eq:unit} (fold-gap $1$) or twice
that (fold-gap $2$).  Both values lie in $(\pi/3,\pi)$ for every
admitted $h\ge 92$: no run of length $\ge 3$ can form from a
constant Chebyshev step, because $r\ge 3$ demands mean interior
step $<\pi/3$~\cite{v2lem}.

\section{The $\gamma$-source formula}\label{sec:gamma}

The folded weights carry an explicit Fej\'er factor
$4\sin^2(\theta/2)=2(1-x)$.  Orthogonal polynomials for the
Jacobi weight $(1-x)$ on $[-1,1]$ (parameters $(\alpha,\beta)=(0,1)$)
have monic recurrence coefficients
\[
\gamma_n=\frac{n(n+1)}{(2n+1)^2},\qquad
\Bigl\lvert\gamma_n-\tfrac14\Bigr\rvert
=\frac{1}{4(2n+1)^2}.
\]
At $n=182$ the right-hand side is $1.9\cdot 10^{-6}$.  The
Fej\'er factor alone does not produce the measured
$O(10^{-2})$ deviation.

\begin{lemma}[Fej\'er Jacobi]\label{lem:fejer}
For the weight $(1-x)$ on $[-1,1]$,
$\lvert\gamma_n-\tfrac14\rvert=1/(4(2n+1)^2)$ exactly
(rational identity).
\end{lemma}

\begin{lemma}[equal-weight cosine grid]\label{lem:eqwt}
On the occupied cosine nodes of FRAME-A $w=9$ ($367$ atoms)
with equal weights, the monic Stieltjes coefficients satisfy
$\gamma_1=\tfrac12$ and $\lvert\gamma_k-\tfrac14\rvert<10^{-5}$
for $2\le k\le 40$ (one fold index of $368$ is unoccupied; the
residual is $10^3$ times smaller than the weighted source).  The node set alone is Chebyshev; the
measured $O(10^{-2})$ is a \emph{weight} perturbation.
\end{lemma}

The source weights are the r342 dictionary (archimedean
digamma $F_A(\xi)=-\log\pi+\Re\psi(\tfrac14+i\xi/2)$ plus the
exact two-cell tent of the prime layer, $v_{\mathrm{pred}}$
relative deviation $<10^{-4}$) times the Fej\'er factor.  The
Jacobi coefficients of $v_2$ are therefore
\[
\gamma_k
=\tfrac14
+\delta_k^{\mathrm{Fejer}}
+\delta_k^{\mathrm{d_{arm}}},
\]
with $\delta^{\mathrm{Fejer}}$ explicit and negligible at
half-filling, and $\delta^{\mathrm{d_{arm}}}$ the only live
term.  This is the source formula.  It is not a bound:
$\delta^{\mathrm{d_{arm}}}$ is the named remainder below.

WKB at the mask $x=\tfrac35$, for a constant $\gamma$ with
$\lvert\gamma-\tfrac14\rvert\le\varepsilon$, the spatial-step
factor against Chebyshev is
$\sqrt{1-x^2}/\sqrt{4\gamma-x^2}$.  At $\varepsilon=\tfrac1{16}$
the WKB band is $[1.332,2.012]$, inside the safe uniform band of
Lemma~\ref{lem:twoperiod}.  The source last-$12$ $\alpha_k$
satisfy $\lvert\alpha\rvert\le 0.071$ (pins), so the turning
points stay near $\pm 1$.

\section{From bounds to $\mathcal{A}_{15}$}\label{sec:a15}

The obstruction of~\cite{v2lem} is a \emph{block}: a long slow
run then three fast steps.  An interleaved two-period
oscillation with the same extremes does not violate.

\begin{lemma}[two-period neighbourhood]\label{lem:twoperiod}
Let $\Delta=\pi/2$ and
$d_i=\Delta\bigl(1+a(-1)^i\bigr)$ for $i=0,\ldots,13$, with
$0\le a\le 0.85$ (all steps in $(0,\pi)$).  Among $2001$
remainders $\eta\in[0,\pi)$, there are $0$ $V_2$-violators.
At $a=0.85$ the step-ratio is $12.3$ and the small step
$0.236<\pi/3$, so $r\ge 3$ dwells are possible: they do not
produce a violator because slow and fast \emph{alternate}.
\end{lemma}

The companion script also gates the blocked rearrangement of
the same values: a slow-then-fast block can leave
$\mathcal{A}_{15}$ at step-ratio $1.83$ (G4 of~\cite{v2lem},
refined).  Ratio is not the obstruction; spatial segregation
is.

\begin{lemma}[dangerous ray]\label{lem:ray}
On $\chi_3$ $w=9$, multiplying the last twelve $\gamma_k$ by a
coherent scale $s$: $s=1.5$ has $0$ $\mathcal{A}_{15}$-violators
($401$ remainders); $s=2.0$ has $2/401$; the actual $v_2$
occupancy first fails at $s=6.5$ (the $8\times$ of~\cite{v2lem}
is past the threshold).  Independent jitter of the last twelve
inside $\lvert\gamma-\tfrac14\rvert\le 0.06$ ($40$ draws) never
leaves $\mathcal{A}_{15}$.
\end{lemma}

The source lives at relative scale $\le 1.24$ about $\tfrac14$
(last-$12$ $\lvert\gamma-\tfrac14\rvert\le 0.052$).  The first
$\mathcal{A}_{15}$ failure on the coherent ray is at scale $2$,
four times farther in the $\varepsilon$-metric
($0.25$ versus $0.06$).  Abstract Jacobi-positivity remains
insufficient (the ray does leave $\mathcal{A}_{15}$); the
source box is far inside the first failure.

\section{The smaller lemma \texorpdfstring{$G_{\varepsilon}$}{Geps}}\label{sec:gbox}

\begin{definition}[$G_\varepsilon$]\label{def:geps}
Let $n=N-2$ be the degree of $v_2$ and let
$(\gamma_k)_{k=1}^{n}$ be the monic Jacobi coefficients of the
bordered chain of $\mu-\nu$.  Condition $G_\varepsilon$ is that
the last twelve coefficients satisfy
\[
\bigl\lvert\gamma_k-\tfrac14\bigr\rvert\le\tfrac1{16}
\qquad\text{and}\qquad
\bigl\lvert\log(\gamma_{k+1}/\gamma_k)\bigr\rvert\le\tfrac25.
\]
\end{definition}

\begin{prop}[reduction]\label{prop:gbox}
$G_\varepsilon$ is strictly more elementary than $V_3'$ in the
sense of the hydra ward of~\cite{v2}~\S5:
\begin{enumerate}
\item \emph{Finite.}  Twelve real inequalities on the chain
  coefficients.  $V_3'$ is a point of $\R^{14}$ together with a
  compact scan of $\eta\in[0,\pi)$.
\item \emph{Closer to the construction.}  The coordinates are
  the three-term recurrence of $\mu-\nu$ (moments, r342
  dictionary, cosine-grid nodes).  No Pr\"ufer phase, no
  Wronskian occupancy, no mask remainder $\eta$.
\item \emph{Fewer quantifiers.}  No $\forall\eta$.
\end{enumerate}
On the $181$-window surface (FRAME-A $89$, FRAME-B $8$,
$\chi_3$ $42$, $\chi_4$ $42$): $G_\varepsilon$ holds, the mask
mesh has $\Delta\theta$-ratio $1$, and the last-$14$ step tuple
lies in $\mathcal{A}_{15}$ ($0$ remainder-violators among $401$
remainders per window; $2001$ on the pins).  Klein-gap windows
kz $51,75,79,103,111$ have last-$3$
$\lvert\gamma-\tfrac14\rvert\le 0.0062$ (smaller than the bulk
small-$N$ worst kz $19$ at $0.0375$).  $\chi_3$ index $16$
produces remainder-$(1,1,1)$ tails, all $V_2$-regular.
The chain is
\[
T_2'
 \;\Longleftarrow\;
 g_i\ge\tfrac38
 \;\Longleftarrow\;
 \mathrm{MED\text{-}CAP}
 \;\Longleftarrow\;
 X_n
 \;\Longleftarrow\;
 V_2
 \;\Longleftarrow\;
 V_3'
 \;\Longleftarrow\;
 G_\varepsilon.
\]
\end{prop}

$G_\varepsilon$ is not a Chebyshev-asymptotic toward L$^*$: it
is a finite box at half-filling, using the cosine-grid identity
rather than Stahl--Totik regularity.  The equal-weight case
(Lemma~\ref{lem:eqwt}) shows the box is a statement about
$d_{\mathrm{arm}}$, i.e.\ about the tent/digamma dictionary,
not about the node set.

\section{Adversaries}\label{sec:adv}

A coherent last-$12$ scale of $2$ leaves $\mathcal{A}_{15}$ for
some $\eta$ (Lemma~\ref{lem:ray}): $V_3'$ uses the source.
Klein-gap windows (r329/r355 families) sit in a \emph{tighter}
coefficient box, not a looser one.  $\chi$ worlds are not
exceptions ($\chi_3$ carries the $v_2$-triple, regular;
$\chi_4$ last-$12$ max $0.0519$, inside $\tfrac1{16}$).
Depth improves the bound (N$=1721$ last-$3$ $0.0011$).
A weight-cluster at the mask still does not fire~\cite{v2lem}.

\section{What remains}\label{sec:rest}

To prove $G_\varepsilon$ cofinally: bound
$\delta_k^{\mathrm{d_{arm}}}$ for $k\in[n-12,n)$ from the
closed-form lag density (r342), uniformly in the window.  The
measured law is that the deviation is largest at small $N$ and
shrinks; the $181$-window surface includes every currently
admitted small-$N$ row.  The implication
$G_\varepsilon+$mesh$\Rightarrow V_3'$ is a Jacobi-on-cosine-grid
statement, gated on the $181$-tuple and on the coherent ray
through scale $1.5$, not an exhaustive $12$-dimensional box
scan.  T1 / $C_K$ unchanged.  No RH claim.

\section{Machine checks}\label{sec:machine}

Every numbered lemma is re-proved in
\texttt{verify\_v3prime.py} (same directory).  Standalone:
$L=4h-2$; Jacobi-$(0,1)$ identity over $\Q$; two-period
$\mathcal{A}_{15}$; contrast regression; WKB band.
Construction: mask $\Delta\theta$-ratio, $G_\varepsilon$ and
Path~A on the pins, equal-weight Chebyshev, the coherent ray,
kz$46$ (step-ratio $5.85$, still in $\mathcal{A}_{15}$).
Optional \texttt{--full}: the $181$-window census.
Exit line: \texttt{V3PRIME LEMMA VERIFIED}.

\section*{Claim boundary}

Finite identities on the cosine grid and the Fej\'er Jacobi
coefficients, a finite two-period neighbourhood of
$\mathcal{A}_{15}$, a named coefficient box $G_\varepsilon$ to
which $V_3'$ reduces, and a $181$-window census.  No statement
about zeros of $\zeta(s)$, in either direction.  No RH claim.

\begin{thebibliography}{9}
\bibitem{v2lem}
S.~Hamann,
\textit{The $V_2$-lemma, reduced},
standalone note, August 2026,
\texttt{rh/problem/v2\_lemma\_v3.tex}.
\bibitem{v2}
S.~Hamann,
\textit{Run-length regularity $V_2$ at the $x$-mask},
standalone note, August 2026,
\texttt{rh/problem/v2\_regularity.tex}.
\bibitem{xn}
S.~Hamann,
\textit{The length invariant $X_n$ of a folded sign field},
standalone note, August 2026,
\texttt{rh/problem/xn\_invariant.tex}.
\end{thebibliography}

\end{document}
